\begin{abstract}
Cross-domain hypothesis generation---discovering mechanistic links between two
domain-specific knowledge graphs---requires traversing graph boundaries that no
single-domain operator can cross.
We present a multi-operator pipeline consisting of \textit{align}, \textit{union},
\textit{compose}, \textit{difference}, and \textit{evaluate} operators, and study its
behaviour on Wikidata-derived bio-chem knowledge graph pairs across three independent
domain contexts (cancer signaling, immunology, and neuroscience).

Our experiments establish three findings.
(1)~The \textit{align} operator is necessary and sufficient to unlock
alignment-dependent cross-domain paths: without it, the compose operator produces
zero cross-domain candidates; with it, the pipeline generates dozens to hundreds of
structurally unique cross-domain path candidates, scaling non-linearly with graph size.
(2)~Raw multi-hop output from real knowledge graphs contains substantial semantic
drift---25\% of deep candidates are structurally non-mechanistic---and a relation-type
filter derived from qualitative review reduces drift to 0\% while preserving all
scientifically meaningful candidates, generalising across all three domain pairs
without parameter retuning.
(3)~The number of unique cross-domain candidates is better predicted by the structural
diversity of bridge entities than by their count: neurotransmitter and eicosanoid bridges
yield 6--8$\times$ more candidates per bridge pair than a single-metabolite bridge.

We further identify a fundamental limitation of the current pipeline: the knowledge
graph encodes relation types at the class level rather than the instance level,
which means filter-passing candidates can carry chemically invalid edges that the
filter cannot detect without substrate-specificity metadata.
This finding, from a post-hoc relation semantics audit (Run~014),
constrains all candidate interpretations and is the central theme of
the Limitations and Discussion sections.
Code and all experimental artefacts will be released upon acceptance
(repository link provided in the supplementary material).
\end{abstract}
